\documentclass[11pt,a4paper,twoside]{article}
\usepackage[utf8]{inputenc}
\usepackage[english]{babel}
\usepackage{actuarialangle, actuarialsymbol}
\usepackage{amsmath,amsfonts,amsthm,amssymb,mathtools}
\begin{document}

% paper size setting in anki \special{papersize=3in,7in}

What are the methods required by law for the actuarial reserve?

For life insurance products, the future benefits are usually supported by preceeding
premium payments. This time difference of cash in/outflow introduces the necessity of an account management.
There are two methods for determining the required reserve
the prospective approach:
present value of future benefits = prospective reserve + present value of future premiums
the retrospective approach:
terminal value of previous premiums = retrospective reserve + terminal value of previous benefits

According to German HGB the prospective approach should always be preferred. The retrospective method
is especially useful under the assumption of stationarity of the distribution (past performance predicts future value).


Give formula's for the actuarial reserve 

Assumptions: \\
Payment times: \( \{0 = s_0, s_1, \dots, s_k \}\) \\
Death benefit: \( D(s_{k+1}) \) if death occurs in \( (s_k, s_{k+1}] \), payable in \( s_{k+1} \) \\
Survival benefit: \( S(s_k) \) if surviving until \( s_k \) , payable in \( s_k \) \\
Premium \( P_x(s_k) \) for \( x \)-year old insured person in period \( (s_k , s_{k+1}] \), payable in advance in \( s_k \) \\
Calculatory maximum age \( w_x \)  \\

Assume \( s_k = k \) then the reserve can be expressed as: 
\[ \Vx[t]{x} = \sum_{l=k}^{w - x} D(l + 1)v^{l + 1 + k} \px[l - k]{x + k} \qx{x + l} + \sum_{l = k}^{w - x} (S(l) - P_x(l)) v^{l - k} \px[l - k]{x + k} \]

For inter-annual time points \( k + u \) with \( u \in [0, 1] \) under the linearity assumption \( \qx[u]{x} = u \qx{x} \)
we get instead
\[ \Vx[k + u]{x} = \frac{v^{1-u}}{1 - u\qx{x+k}} (D(k+1)(1 - u)\qx{x+k} + \Vx[k+1]{x} \px{x+k} \]

Example \( n \)-year insurance for survival with insured sum \( 1 \) and constant annual premiums
\( P_{x,n}(\Ex[n]{x}) \) due in advance:
\[ \Vx[k]{x}(\Ex[n]{x}) = \Ex[n-k]{x+k} - P_{x,n}(\Ex[n]{x})\ax**{x+k: \angl n - k} = \frac{D_{x+n}}{D_{x+k}}(1 - \frac{N_{x+k} - N_{x+n}}{N_x - N_{x+n}}) \]

The actuarial balance equation or fundamental equation is a recursive expression for \( \Vx[k]{x} \) which allows to fast track some computations: \\
Forward recursion (prospective):
\[ \Vx[k+1]{x} = \frac{1}{v \px{x+k}}(\Vx[k]{x} - S(k) + P_x(k) - v\qx{x+k} D(k+1)) \]
with \( \Vx[0]{x} = 0 \). \\
Backward recursion (retrospective):
\[ \Vx[k]{x} = S(k) - P_x(k) + v \qx{x+k} D(k+1) + v\px{x+k} \Vx[k+1]{x} \]
with \( \Vx[n]{x} = S(n) \) or \( 0 \) if no final survival benefit is included. \\

The premia can be decomposed into a savings and risk component which are expressed in terms of the reserve:
\[ P^S_x(k) = \frac{K(k)}{K(k+1)} \Vx[k+1]{x} - \Vx[k]{x} \]
The savings component is for the adjustment of the prospetive reserve at the end of the year.
\[ P^R_x(k) = \frac{K(k)}{K(k+1)}(D(k+1) - \Vx[k+1]{x} \mathbb{P}[T_x \leq k+1 | T_x > k]  ) \]
The risk component is for covering the risk (of death) in the period \( (k, k+1] \) and it corresponds
to the premium of a term life insurance in \( (k, k+1] \) with insured sum \( D(k+1) - \Vx[k+1]{x} \).
\end{document}
